\documentclass[10pt,a4paper]{article}
\usepackage{amsmath,amssymb,amsthm,graphicx}
\usepackage{hyperref}
\usepackage{latexsym}
\usepackage{amssymb}
\usepackage{amsthm}
\usepackage{ulem}
\usepackage{enumitem}
\setlist{itemsep=0.4pt,topsep=0pt}
\parindent 0pt
\parindent 0pt
\parskip 12pt
\newcommand{\HCF}{\mathrm{hcf}}
\newcommand{\ZZ}{\mathbb{Z}}
\newcommand{\rk}{\mathrm{rk}}
\DeclareMathOperator{\len}{length}
\DeclareMathOperator{\bl}{b}
\DeclareMathOperator{\cl}{cl}
\DeclareMathOperator{\di}{d}
\DeclareMathOperator{\Cl}{Cl}
\DeclareMathOperator{\acl}{acl}
\DeclareMathOperator{\tp}{tp}
\DeclareMathOperator{\Th}{Th}
\newcommand{\Cor}{\mathcal{C}_\mathrm{or}}
\newcommand{\Cun}{\mathcal{C}_\mathrm{un}}
\newcommand{\mons}{\mathbb{M}}
\newcommand{\strong}{\mathrm{strong}}
\newcommand{\axioms}{\Delta_{\mathrm{FSTS}}}
\newtheoremstyle{note}% hnamei
{8pt}% hSpace abovei
{14pt}% hSpace belowi
{\normalfont}% hBody fonti
{}% hIndent amounti1
{\bfseries}% hTheorem head fonti
{ }% hPunctuation after theorem headi
{.5em}% hSpace after theorem headi2
{}% hTheorem head spec (can be left empty, meaning ‘normal’)i
\theoremstyle{note}
\newtheorem{thm}{Theorem}[section]
%\newtheorem{theorem}{Theorem}
\newtheorem{lemma}[thm]{Lemma} 
\newtheorem{df}[thm]{Definition}
\newtheorem{rmk}[thm]{Remark} 
\newtheorem{prop}[thm]{Proposition}
\newtheorem{fact}[thm]{Fact}
\newtheorem{cor}[thm]{Corollary}
\newtheorem{conj}[thm]{Conjecture}
\newtheorem{ex}[thm]{Example}
\newtheorem{qu}[thm]{Question}
\newtheorem*{notation}{Notation} 
\title{The generic structure is the free Steiner triple system on countably many generators}
\author{Silvia Barbina, Enrique Casanovas}
\date{arxiv:2411.13723v3; Creative Commons Attribution 4.0 International}
\begin{document}
\maketitle
\noindent\textit{Editorial note.} The counted statement is the author's environment at source0.tex lines 760-761 and its proof environment opens at line 762; the pairing is by adjacency and the retained proof range is the whole environment.\par\medskip
\noindent\textit{Editorial note.} The preamble is the author's own package lines, macro block and theorem-environment declarations (source lines 4-65); the author's optioned amsmath line and the fleqn-dependent layout line are not retained.\par\medskip
\noindent\textit{Editorial note.} The source carries no label-number dump and no counter values are imposed, so the printed numbering is the wrapper's default; identity is the author label, or the pinned source position where the author gave none.\par\medskip
\noindent\textit{Editorial note.} Metadata provenance: the source record of the pinned copy and the licence evidence file are both present; the licence is CC BY 4.0 with status supported, and the card's source\_version matches the value recorded in that evidence file.\par\medskip
\noindent\textit{Editorial scope.} One original result of the pinned author manuscript, counted once: the generic is the free system on countably many generators (author label Feb\_23\_2024\_4).\par\medskip
\noindent\textit{Allowed standard prerequisites (editorial).}\begin{itemize}
\item The generic structure. The generic is the model obtained by the amalgamation construction from the finite unconfined systems, with a distinguished set of vertices.. It is the structure shown to be free.
\item Freely generated system. A system is freely generated by a set when that set is independent and the system is generated by it.. It is the conclusion of the statement.
\item Extension of finite embeddings. Finiteness of the generating data makes every partial embedding from a finite subset extend to an embedding of the generated subsystem.. It is what forces the universal mapping property.
\end{itemize}\medskip
	\begin{prop}\label{May12_2.2}Assume $M$ is a Steiner triple system and $A\subseteq M$. The following are equivalent:
		\begin{enumerate}
			\item $A$ is a free base of $M$ (as a Steiner quasigroup)
			\item There is a well-founded F-ordering of $M$ over $A$
			\item There is a well-founded HF-ordering of $M$ and $A$ is the set of points that are not a product of smaller elements.
		\end{enumerate}
		\end{prop}

	\begin{prop}\label{Feb_23_2024_4} The generic is isomorphic to the free Steiner triple system with $\omega$ generators.
	\end{prop}

	\begin{proof} First note that the generic $F$ is an STS, that is, the product of two distinct elements $a_1,a_2$ is defined: let $b$ be a new vertex and let $A=\{a_1,a_2,b\}$ be an STS (so $ \{a_1, a_2, b \}$ is a block). Then $A\in\Cun$.  Let $B \leq F$ be such that $  \{a_1, a_2\} \subseteq B$. If $a_1 \cdot a_2 \notin B$, then $a_1 \cdot a_2$ can be put on top in an HF-ordering of $B$ and, therefore, $B \leq B(a_1 \cdot a_2)$. So there is an embedding $f:B(a_1 \cdot a_2)\rightarrow F$ over $B \supseteq\{a_1,a_2\}$. Moreover, $F$ must be a free Steiner triple system, since it has a well-founded HF-ordering.  On the other hand, the generic can be constructed by  including at each stage an element which does not form a block with any previous elements. Then the corresponding  HF-ordering has infinitely many elements which do not form a block with any preceding elements. By Proposition~\ref{May12_2.2}, these elements belong to a free base of~$F$.
	
%	 Then $A\in\Cun$ and $\{a_1,a_2\}\leq A$, and, therefore, there is an embedding $f:A\rightarrow F$ over $\{a_1,a_2\}$. Moreover, $F$ must be a free Steiner triple system, since it has a well-founded HF-ordering.  On the other hand, the generic can be constructed by  including at each stage an element which does not form a block with any previous elements. Then the corresponding  HF-ordering has infinitely many elements which do not form a block with any preceding elements. By Proposition~\ref{May12_2.2}, these elements belong to a free base of~$F$.
	\end{proof}

\end{document}
